% !TeX program = pdflatex
\documentclass[11pt,a4paper]{article}
\usepackage[T1]{fontenc}
\usepackage[utf8]{inputenc}
\usepackage[polish]{babel}
\usepackage{lmodern}
\usepackage{amsmath,amssymb,amsthm,mathtools}
\usepackage{microtype}
\usepackage{xcolor}
\usepackage{geometry}
\usepackage{enumitem}
\usepackage{fancyhdr}
\usepackage[unicode]{hyperref}
\hypersetup{colorlinks=true,linkcolor=blue!45!black,urlcolor=blue!45!black}
\geometry{margin=2.5cm}
\setlength{\parindent}{0pt}
\setlength{\parskip}{4pt}
\setlength{\emergencystretch}{2em}
\setlist[enumerate]{topsep=4pt,itemsep=2pt}
\setlist[itemize]{topsep=4pt,itemsep=2pt}
\numberwithin{equation}{section}
\theoremstyle{plain}
\newtheorem{twierdzenie}{Twierdzenie}[section]
\newtheorem{lemat}[twierdzenie]{Lemat}
\newtheorem{stwierdzenie}[twierdzenie]{Stwierdzenie}
\newtheorem{wniosek}[twierdzenie]{Wniosek}
\theoremstyle{definition}
\newtheorem{definicja}[twierdzenie]{Definicja}
\newtheorem{przyklad}[twierdzenie]{Przykład}
\theoremstyle{remark}
\newtheorem{uwaga}[twierdzenie]{Uwaga}
\renewcommand{\proofname}{Dowód}
\pagestyle{fancy}
\fancyhf{}
\fancyhead[L]{\small Rozszerzenia pojęcia całki}
\fancyhead[R]{\small\ifnum\value{section}>0 Sekcja \thesection\else Wprowadzenie\fi}
\fancyfoot[C]{\thepage}
\renewcommand{\headrulewidth}{0.4pt}
\newcommand{\R}{\mathbb R}
\newcommand{\N}{\mathbb N}
\newcommand{\C}{\mathbb C}
\newcommand{\Q}{\mathbb Q}
\newcommand{\La}{\lambda}
\newcommand{\ind}[1]{\mathbf 1_{#1}}
\newcommand{\osc}{\operatorname{osc}}
\newcommand{\Var}{\operatorname{Var}}
\newcommand{\mesh}{\operatorname{mesh}}
\newcommand{\supp}{\operatorname{supp}}
\newcommand{\AC}{\mathrm{AC}}
\newcommand{\BV}{\mathrm{BV}}
\newcommand{\HK}{\mathrm{HK}}
\newcommand{\RS}{\mathrm{RS}}
\newcommand{\KS}{\mathrm{KS}}
\newcommand{\Dn}{\underline D}
\newcommand{\Du}{\overline D}
\newcommand{\pv}{\mathrm{p.w.}}
\DeclareMathOperator{\esssup}{ess\,sup}
\title{Rozszerzenia pojęcia całki\\\large Definicje, kryteria, dowody i przykłady}
\author{}
\date{}